\chapter{货币、支付与流动性架构}
\label{chap:Money_Payments_and_Liquidity_Architecture}

\begin{openingcase}{支付应用并不是结算资产}

到 2010 年代末，移动支付和网络支付已深深融入中国的日常商业活动。对用户而言，交易往往十分简单：打开应用、扫码、确认金额，支付就完成了。然而，在这个界面背后，支付的制度路径正在变化。

此前，非银行支付机构直接连接多家商业银行，处理涉及银行账户的网络支付。中国人民银行要求支付机构切断这些直连，并通过集中清算基础设施处理相关交易。截至 2019 年 1 月 14 日，非银行支付机构与商业银行共同开展的支付业务已通过网联清算或中国银联处理。\footnote{中国人民银行，《中国人民银行年报 2018》，有关网联平台及非银行支付机构从与商业银行直连模式迁移的讨论。}

对大多数客户而言，支付体验基本没有变化，尽管其背后的制度架构已发生重大改变。

\end{openingcase}

%========================================================
% Core Questions
%========================================================

\begin{tcolorbox}[
  enhanced,
  colback=FinanceBlue!6,
  colframe=FinanceBlue,
  boxrule=0.6pt,
  sharp corners,
  left=8pt,
  right=8pt,
  top=7pt,
  bottom=7pt,
  before skip=1em,
  after skip=1.6em,
  title=\textbf{核心问题},
  coltitle=white,
  colbacktitle=FinanceBlue,
  fonttitle=\bfseries,
  attach boxed title to top left={xshift=6pt,yshift=-2mm},
  boxed title style={
    boxrule=0pt,
    sharp corners,
    left=5pt,
    right=5pt,
    top=2pt,
    bottom=2pt
  }
]
\begin{enumerate}[
  leftmargin=*,
  label=\textcolor{FinanceBlue}{\textbf{\arabic*.}},
  itemsep=0.55em,
  topsep=0.2em
]
\item 为什么有些金融请求权能够充当货币，而另一些不能？
\item 支付指令如何经过清算和结算处理？支付何时具有最终性？
\item 为什么现代货币体系同时依赖中央银行货币和商业银行货币？什么机制支持二者按面值兑换？
\end{enumerate}
\end{tcolorbox}

\section{作为金融约定的货币}

为什么一个人愿意接受他人提供的某种物品或账户余额，以交换商品、劳动或金融请求权？答案不仅取决于货币的形式，还取决于使货币被广泛接受、并能用于清偿义务的制度。

\subsection{作为社会与制度安排的货币}

货币的有用性，依赖于人们共同预期他人也会接受它作为支付。如果每个卖方在接受支付之前，都必须检查发行人的资产支持、法律地位和资产的可接受性，日常交换就会变得昂贵。货币通过建立普遍接受的计量单位和结算惯例，减轻了这种负担。

这一特性具有一定的自我强化作用。接受某种资产的人越多，它用于支付的价值就越高；接受范围越广，其他人持有和使用它的理由也越充分。这种反馈形成货币交换中的网络效应（network effects）。但仅有网络效应还不够，参与者还必须相信货币单位将持续可用，以其进行的支付会被认可并得到结算。
现代货币体系通过法律、支付基础设施、中央银行安排、监管和共同预期来支持这种信心。

货币并不只是一种社会惯例。其使用既依赖共同接受，也依赖支持货币请求权的制度安排。例如，银行存款是存款人的资产，也是发行银行的负债。因此，货币的普遍可接受性既建立在共同预期上，也建立在支持支付和结算的安排上。

\subsection{货币的核心职能}

一项货币请求权在得到广泛接受和制度支持后，就能够履行三项核心货币职能。

\emph{记账单位}（unit of account）为价格和金融义务提供共同语言。贷款协议可以规定支付人民币 100 万元，而不必约定一组不断变化的商品。债券可以用与记录发行人收入、投资者财富相同的单位，承诺支付利息和本金。缺乏共同单位，会增加比较和缔约的成本。

\emph{交换媒介}（medium of exchange）职能使货币能够促进交易并清偿义务。现金可以直接结清面对面的交易，银行存款则可以通过支付系统转移。


\begin{table}[H]
\centering
\caption{货币的核心职能}
\label{tab:ch2-money-functions}
\renewcommand{\arraystretch}{1.4}
\begin{tabularx}{\textwidth}{
    @{}
    >{\raggedright\arraybackslash}p{3.0cm}
    >{\raggedright\arraybackslash}X
    >{\raggedright\arraybackslash}X
    @{}
}
\toprule
\textbf{职能}
& \textbf{经济作用}
& \textbf{金融意义} \\
\midrule

记账单位
& 为价格与义务提供共同计量标准
& 使契约、会计核算和估值具有可比性 \\

交换媒介
& 便利交易并清偿义务
& 支持商品和服务的高效交换 \\

价值储藏
& 将购买力跨期保存
& 使财富以具有流动性的名义形式持有 \\

\bottomrule
\end{tabularx}
\end{table}

\emph{价值储藏}（store of value）职能使购买力能够跨期保存。货币并非完美的价值储藏手段，因为通货膨胀会削弱其实际购买力。即便如此，它通常仍比许多其他资产更稳定、流动性更高。

这些职能为区分货币、其他金融资产与支付工具提供了有用依据。例如，国债可以是良好的价值储藏手段，却不充当交换媒介。支付卡便利交易，但本身并不是货币：借记卡使持卡人能够使用银行存款，信用卡则使其能够动用授信额度。一种外币在某一司法辖区可能履行货币职能，在另一辖区则主要作为金融资产持有。


\subsection{货币的资产负债表结构}

现代经济中的大部分货币，表现为记录在机构资产负债表上的金融请求权。最重要的形式是商业银行存款、中央银行准备金和通货，尽管它们在持有资格和用途上有所不同。

商业银行存款是银行的负债、存款人的资产。中央银行准备金是中央银行的负债、合资格账户持有人的资产，这些持有人主要是银行和其他指定机构。通货也记录为中央银行负债，尽管现代法定货币通常不能按固定比率兑换为黄金或另一种指定资产。

\begin{table}[H]
\centering
\caption{银行存款同时是资产和负债}
\label{tab:ch2-deposit-taccount}

\renewcommand{\arraystretch}{1.25}

\begin{tabularx}{0.9\textwidth}{
    @{}
    >{\raggedright\arraybackslash}X
    >{\raggedright\arraybackslash}X
    >{\raggedright\arraybackslash}X
    @{}
}
\toprule
\multicolumn{1}{c}{\textbf{家庭}}
&
\multicolumn{2}{c}{\textbf{商业银行}}
\\
\cmidrule(r){1-1}
\cmidrule(l){2-3}

\textbf{资产}
&
\textbf{资产}
&
\textbf{负债}
\\

银行存款：人民币 100 元
&
--
&
银行存款：人民币 100 元
\\

\bottomrule
\end{tabularx}
\end{table}

表~\ref{tab:ch2-deposit-taccount} 展示了同一笔人民币 100 元存款如何出现在两张资产负债表上：它是家庭的资产，同时是银行的负债。

资产负债表视角说明，不同形式的货币在经济上并不完全相同。
人民币 100 元银行存款和 100 元纸币采用相同记账单位，通常按面值被接受，但它们的发行人和资产负债表关系不同。存款是商业银行的负债、存款人的资产；纸币则是中央银行的负债、持有人的资产。

\begin{table}[H]
\centering
\caption{部分货币形式及其资产负债表地位}
\label{tab:ch2-money-claims}

\renewcommand{\arraystretch}{1.15}
\setlength{\tabcolsep}{7pt}

\begin{tabularx}{\textwidth}{
    @{}
    >{\raggedright\arraybackslash}p{3.35cm}
    >{\raggedright\arraybackslash}p{2.85cm}
    >{\raggedright\arraybackslash}X
    @{}
}
\toprule
\textbf{货币形式}
& \textbf{发行人}
& \textbf{资产负债表关系} \\
\midrule

通货\newline
{\small （纸币与硬币）}
& 中央银行
& 中央银行负债；公众持有的资产 \\
\addlinespace[0.3em]

中央银行准备金
& 中央银行
& 中央银行负债；合资格账户持有人的资产 \\
\addlinespace[0.3em]

商业银行存款
& 商业银行
& 银行负债；存款人的资产 \\

\bottomrule
\end{tabularx}
\end{table}

同样的资产负债表逻辑，也有助于说明商业银行存款如何创造。假设银行新发放人民币 100 元贷款，并在借款人的存款账户中记入同额资金。在银行的资产负债表上，

\begin{equation}
\Delta \text{贷款}=+100,
\qquad
\Delta \text{存款}=+100.
\label{eq:ch2-loan-deposit}
\end{equation}

因此，贷款在发放时创造了一笔新存款。借款人同时获得存款资产和相应的贷款负债，其净财富并未等额增加。存款创造仍然受到银行整体资产负债表状况和监管约束的限制。

\subsection{信任、可信度与货币稳定}

货币降低信息成本，因为人们不必在每笔交易中重新评估每项货币请求权背后的发行人。银行转账的收款人通常依赖存款的可转移性，而不必评估付款人的开户银行；通货持有人也不必检查中央银行的资产负债表。货币交换的效率，建立在无需反复核实便能将被接受的货币请求权视为可靠的基础上。

这种信任来自支持货币体系的制度和安排的可信度。货币单位必须足够稳定，才能支持定价和缔约；货币请求权必须能够按承诺转移或兑换。可靠的支付系统，以及明确所有权和结算最终性的法律规则，也支持这种可信度。

货币信任下降会通过多种方式影响行为。如果存款人担心银行无法满足提款需求，可能将存款换成通货，或将资金转至其他银行。如果对私人发行货币请求权的信心减弱，该请求权可能折价交易，或不再被接受。如果通胀削弱人们对记账单位的信心，家庭和企业可能缩短契约期限、对价格实行指数化，或转向其他货币与资产。

货币稳定还取决于这样的预期：以同一单位计价的不同货币形式，将保持可互换性。这一预期支持整个经济中的共同定价、缔约和支付。

\section{支付系统与金融网络}

货币只有能够在用户之间转移，才能充当支付手段。支付系统（payment system）是用于转移资金、清偿货币义务的一整套工具、规则、机构、技术与结算安排。

\subsection{支付工具与传统支付基础设施}

支付工具用于发起和传递支付指令。现金将支付工具与结算资产合为一体：交付通货即可直接完成交易。大多数非现金支付，则将支付指令与用于结算的资产分开。

非现金支付工具在由谁发起支付、如何传递指令方面存在差异。支票和贷记转账由付款人发起，直接借记则由收款人在事先授权下发起。银行卡和移动支付依靠支付网络及数字界面来认证用户、授权交易和传送支付指令。

\begin{figure}[H]
\centering

\resizebox{0.94\textwidth}{!}{%
\begin{tikzpicture}[
  >=Latex,
  box/.style={
  draw,
  rounded corners,
  align=center,
  minimum height=0.85cm,
  text width=1.5cm,
  inner xsep=2pt,
  inner ysep=3pt,
  line width=0.8pt
  },
  flow/.style={
    -Latex,
    line width=0.85pt,
    FinanceBlue
  },
  note/.style={
    font=\scriptsize,
    align=center,
    text=FinanceBlue,
    fill=white,
    inner sep=1pt
  }
]

\node[box] (payer) at (0,0) {付款人};
\node[box] (bankA) at (3.6,0) {付款人银行};
\node[box, text width=2.2cm] (infra) at (6.3,0)
{清算与\\结算系统};

\node[box] (bankB) at (9.0,0) {收款人银行};
\node[box] (payee) at (11.8,0) {收款人};

\draw[flow] (payer) -- 
  node[midway, above=6pt, note]{支付\\指令}
  (bankA);

\draw[flow] (bankA) -- (infra);

\draw[flow] (infra) -- (bankB);

\draw[flow] (bankB) -- 
  node[midway, above=5pt, note]{贷记}
  (payee);

\end{tikzpicture}%
}

\caption{简化的跨行支付链条}
\label{fig:ch2-payment-chain}
\end{figure}

图~\ref{fig:ch2-payment-chain} 展示了一条简化的跨行支付链条。如果付款人和收款人在同一家银行开户，交易可以通过减少一个存款账户、增加另一个账户的余额在行内结算；若双方使用不同银行，支付就会产生银行间义务。参与银行随后必须相互结算，通常使用双方均接受的结算资产。

因此，支付网络不只是一套通信系统。其设计影响义务何时产生和清偿、参与者需要多少流动性，以及结算最终完成之前面临的信用和操作风险敞口。

\subsection{清算、轧差与支付义务}

清算（clearing）是传送、核对和确认支付指令，并将其转化为待结算义务的过程。在适用的情况下，这些义务可在结算前进行轧差。

考虑一组分别以 $i$ 和 $j$ 标识的银行。令 $p_{ij}$ 表示银行 $i$ 必须向银行 $j$ 支付的金额，则银行 $i$ 的付款总额为

\begin{equation}
G_i^{\mathrm{out}}
=\sum_{j\neq i} p_{ij},
\label{eq:ch2-gross-outgoing}
\end{equation}
其净结算头寸可表示为

\begin{equation}
N_i=\sum_{j\neq i} p_{ij}-\sum_{j\neq i} p_{ji}.
\label{eq:ch2-net-position}
\end{equation}

$N_i$ 为正表示净付款义务，为负表示净收款。在封闭的清算系统中，

\begin{equation}
\sum_i N_i=0.
\label{eq:ch2-net-zero}
\end{equation}

轧差（netting）可以显著降低结算所需流动性。假设银行 A 欠银行 B 100，银行 B 欠银行 C 90，银行 C 欠银行 A 80。表~\ref{tab:ch2-netting-example} 列示了相应的总额流量和净结算头寸。

\begin{table}[H]
\centering
\caption{总额支付与净结算头寸}
\label{tab:ch2-netting-example}
\begin{tabular}{@{}lrrr@{}}
\toprule
\textbf{银行} & \textbf{付款} & \textbf{收款} & \textbf{净头寸} \\
\midrule
A & 100 & 80 & 支付 20 \\
B & 90 & 100 & 收到 10 \\
C & 80 & 90 & 收到 10 \\
\bottomrule
\end{tabular}
\end{table}

三笔总额支付合计为 270，但经过多边轧差后，系统只需结算合计 20 的净借方义务。银行 A 支付 20，银行 B 和 C 各收到 10。轧差降低了结算所需流动性。

在延迟净额结算系统中，节省流动性也有代价。支付义务可能直到清算周期结束才结算，使参与者面临其他参与者违约的风险。支付系统设计需要在流动性效率、结算最终性和风险控制之间权衡。

\subsection{结算、支付最终性与结算风险}

结算（settlement）是通过结算机构账簿上的转账来清偿支付义务。支付最终性（payment finality）是指根据适用的系统规则和法律，支付被视为完成且不可撤销的时点。

一笔简单的跨行转账可以说明清算与结算的区别。假设银行 A 的客户向银行 B 的客户支付人民币 100 元。如果支付用中央银行准备金结算，主要的资产负债表变化如表~\ref{tab:ch2-cross-bank-payment} 所示。

\begin{table}[H]
\centering
\caption{人民币 100 元跨行支付的资产负债表影响}
\label{tab:ch2-cross-bank-payment}

\renewcommand{\arraystretch}{1.2}

\begin{tabularx}{\textwidth}{
    @{}
    >{\raggedright\arraybackslash}p{3.0cm}
    >{\raggedright\arraybackslash}X
    >{\raggedright\arraybackslash}X
    @{}
}
\toprule
\textbf{机构}
& \textbf{资产}
& \textbf{负债} \\
\midrule

银行 A
& 准备金 $-100$
& 付款人存款 $-100$ \\

银行 B
& 准备金 $+100$
& 收款人存款 $+100$ \\

中央银行
& ---
& 对银行 A 的准备金负债 $-100$；对银行 B 的准备金负债 $+100$ \\

\bottomrule
\end{tabularx}
\end{table}

准备金从银行 A 转至银行 B，清偿了银行间义务。客户的存款并非从一家银行“搬运”到另一家银行，而是银行 A 减少对付款人的负债，银行 B 增加对收款人的负债，同时中央银行准备金在两家银行之间转移。当这一结算根据系统规则成为不可撤销、无条件的转移时，就达到了结算最终性。

结算资产应使义务得以清偿，而不在结算时引入显著的额外信用风险。中央银行货币被广泛用于具有系统重要性的国内支付系统结算。

在达到结算最终性之前，参与者仍面临结算风险：预期结算可能延迟，或无法完成。这类失败可能引发几种相关风险，如表~\ref{tab:ch2-payment-risk} 所示。

\begin{table}[H]
\centering
\caption{支付系统风险的主要形式}
\label{tab:ch2-payment-risk}

\renewcommand{\arraystretch}{1.2}

\begin{tabularx}{\textwidth}{
    @{}
    >{\raggedright\arraybackslash}p{3.0cm}
    >{\raggedright\arraybackslash}X
    @{}
}
\toprule
\textbf{风险}
& \textbf{经济含义} \\
\midrule


信用风险
& 参与者可能在到期时或此后任何时点，都无法全额履行义务。 \\

流动性风险
& 参与者可能无法在到期时履行义务，但可以在较晚时间全额结算。 \\

操作风险
& 技术故障、网络事件、人为差错或基础设施中断，可能妨碍处理或结算。 \\

法律风险
& 规则、所有权、轧差或最终性方面的不确定性，可能削弱可执行性。 \\

\bottomrule
\end{tabularx}
\end{table}

这些风险的重要性，部分取决于结算的组织方式。实时全额结算（real-time gross settlement）系统逐笔、连续地结算支付，缩短银行间风险敞口的持续时间，却增加日间流动性需求。延迟净额结算（deferred net settlement）系统则先累积义务、进行轧差，再结算，因而降低流动性需求，却使风险敞口持续更久。

结算设计需要在流动性效率、实现最终性的速度和风险控制之间权衡。

\subsection{数字支付网络与平台}

数字化改变支付的发起、认证和路由方式，但没有消除区分货币请求权、清算与结算的必要。移动钱包或支付应用可能让用户感觉支付即时完成，但其底层价值仍可能是银行存款、受监管的支付请求权或另一种货币资产。关键问题不仅是如何支付，还包括转移了什么请求权、如何完成最终结算。

支付网络连接用户、商户、金融机构、支付服务提供商，以及清算或结算基础设施。其价值不仅依赖技术，还依赖参与范围、互操作性、规则和信任。更大的网络可以使支付更方便、更便宜，但也可能造成对少数平台或基础设施的依赖。

\begin{figure}[H]
\centering

\begin{tikzpicture}[
  >=Latex,
  layer/.style={
    draw,
    rounded corners,
    align=center,
    minimum height=0.95cm,
    text width=0.82\textwidth,
    inner xsep=6pt,
    inner ysep=5pt,
    line width=0.7pt,
    font=\footnotesize
  },
  arrow/.style={
    -Latex,
    line width=0.75pt,
    FinanceBlue
  }
]

\node[layer] (interface) at (0,2.4)
{\textbf{用户界面层}\\
移动钱包、支付应用、二维码、银行卡或在线收银台};

\node[layer] (network) at (0,0.7)
{\textbf{支付网络与平台层}\\
用户、商户、金融机构、支付服务提供商、规则、互操作性与信任};

\node[layer] (claim) at (0,-1.0)
{\textbf{货币请求权层}\\
银行存款、受监管的支付请求权或其他货币资产};

\node[layer] (settlement) at (0,-2.7)
{\textbf{清算与结算层}\\
计算义务、在适用时轧差，并以结算资产完成最终结算};

\draw[arrow] (interface) -- (network);
\draw[arrow] (network) -- (claim);
\draw[arrow] (claim) -- (settlement);

\end{tikzpicture}

\caption{数字支付的界面、网络、请求权和结算层次}
\label{fig:ch2-digital-payment-layers}
\end{figure}

图~\ref{fig:ch2-digital-payment-layers} 表明，数字支付涉及的并不只是用户界面的变化。同样的可见支付行为，可能基于不同的底层货币请求权、网络安排和结算资产。

在这种分层架构中，平台层协调用户、商户、银行和支付服务提供商之间的互动。平台通过接入管理、身份认证、交易路由和数据处理，降低交易成本与处理延迟。这些安排也产生网络效应：支付服务对参与者的价值可能随着用户和商户数量增加而提高，在创造规模经济的同时，也可能强化市场集中。支付数据能够改善欺诈识别和服务设计，但其收集与使用也引发隐私、市场势力、互操作性和接入方面的问题。

同样的逻辑适用于较新的支付技术。
可编程转账和代币化请求权，可以通过在交易逻辑中嵌入条件，使交易、清算和结算更加紧密地结合。人工智能能够改善欺诈识别、交易监测、流动性预测和支付路由。然而，更快的执行和更好的预测，本身并不能确立法律最终性、保证按面值赎回或消除流动性需求。这一经济问题既涉及计算，也涉及制度。

\section{货币体系中的流动性}

只有在义务到期时，货币请求权能够被使用、转移或兑换，货币才能充当支付手段。因此，支付系统不仅需要货币资产，还需要流动性，即在需要时获得所需结算资产的能力。

\subsection{流动性、可兑换性与经济交换}

流动性不同于资产价值，它描述资产能够多容易地被用于履行到期支付或结算义务。通货可以直接清偿许多零售支付。政府证券可能价值很高、交易活跃，但通常必须先通过出售或抵押借款，将其价值转化为可支用资金，才能支持普通支付。区别在于持有价值与完成支付之间的距离。

可兑换性（convertibility）能够缩短这一距离。这里，可兑换性指将一种货币请求权兑换为另一种的能力，通常按面值且没有显著延迟。银行存款是商业银行的私人负债，但当它能够按面值兑换为通货，并通过支付网络转移时，就能在日常交换中发挥货币作用。

当按面值赎回的信心减弱时，可兑换性的重要性就会显现。存款人可能提取现金，或将余额转至其他机构。更一般地说，不能可靠地按其标示价值赎回的支付请求权，可能不再作为货币的紧密替代物流通。

不过，可兑换性并不消除时间的重要性。假设一家银行预计下午 4 点有大笔资金流入，却必须在上午 10 点付款。即使日终头寸可以平衡，它仍需要上午 10 点的流动性。这种时间错配产生了对准备金、日间信贷，以及其他能在需要时提供结算资产的安排的需求。


\begin{tcolorbox}[
  enhanced,
  breakable,
  colback=white,
  colframe=FinanceBlue,
  boxrule=0.7pt,
  sharp corners,
  borderline west={2.2pt}{0pt}{FinanceBlue},
  title=\textbf{金融视角：可兑换性与流动性},
  coltitle=white,
  colbacktitle=FinanceBlue,
  fonttitle=\bfseries,
  before skip=1.2em,
  after skip=1.2em
]
可兑换性关注金融请求权能否兑换为可用于支付或结算的资产，通常按面值兑换。流动性还要求相关资产在需要的时间、以需要的数量可供使用。因此，一项请求权原则上可兑换，却未必能在义务到期时提供足够流动性。
\end{tcolorbox}

在系统层面，同样的时间错配会引发支付梗阻。机构可能等待收款后才对外付款。当这种行为普遍出现时，即使全日总流动性足够，结算仍可能延迟。因此，流动性安排不仅必须保证结算资产数量充足，还必须保证资产在参与者之间及时可用。

维持流动性具有机会成本。准备金和其他高流动性资产提供即时结算能力，但收益可能低于流动性较低的替代资产。因此，金融机构需要在流动性与收益之间权衡：一方面持有足以履行义务的流动资产，另一方面将其他资金配置到收益更高的用途。

这种权衡反映了流动性和可兑换性在经济交换中的更广泛作用。它们减少支付前变现其他资产的需要，缩短义务清偿的延迟，从而降低交易摩擦，支持货币请求权在整个经济中流通。

\subsection{安全资产与流动性需求}

流动性需求有助于解释安全资产在货币体系中的特殊作用。安全且具有流动性的资产，能够被调动以满足支付和结算需求，同时限制信用损失敞口。

安全性与流动性在概念上不同。这里，安全性主要指信用质量：安全资产承诺的支付无法按期足额实现的风险较低。流动性则关注资产能否在支付或结算义务到期时被便捷地使用或调动。

仅有安全性并不决定资产的流动性。流动性还取决于能否以有利条件迅速出售或质押资产，尤其是在市场承压时。

流动性和抵押服务也会反映在资产收益率中。在其他条件可比的资产中，投资者可能愿意接受较低收益率，以换取更有价值的流动性或抵押服务。这些服务所带来的价值，通常称为便利收益（convenience yield）。

\begin{equation}
y
=
y^{*}-\lambda,
\qquad \lambda>0,
\label{eq:ch2-convenience-yield}
\end{equation}
其中，$y^{*}$ 是其他方面可比、但不提供同等便利服务的资产收益率，$\lambda$ 表示与流动性、抵押及相关服务有关的便利收益。该表达式是一种经济分解，而非定价模型；$\lambda$ 可能随市场条件变化。

在不确定性上升或金融压力时期，对安全和流动资产的需求通常增加。市场参与者可能减少风险较高请求权的持有，转向现金、短期政府证券，或对合资格机构而言，转向中央银行准备金。这种“逃向安全”（flight to safety）的行为能够降低个体风险敞口，却可能收紧其他领域的融资条件和市场流动性。

在机构层面，保全流动性的努力可能放大压力。如果机构同时减少贷款、保留结算资产或出售流动性较差的证券，融资条件可能收紧，资产市场流动性可能恶化。对单个机构审慎的行动，也可能促成全系统流动性压力。

当这种行为广泛出现时，流动性需求就成为整个系统的协调问题。有韧性的货币体系需要在条件变化时支持流动性调动和分配的安排。中央银行流动性工具、抵押品框架和支付系统安排，都是这一制度架构的组成部分。


\section{公共货币与私人货币的架构}
现代货币体系由相互联系的公共和私人资产负债表所发行的货币负债组成。中央银行货币提供共同的结算基础，而商业银行存款构成日常经济活动中所用货币的很大一部分。

\subsection{公共货币与中央银行}

现代货币体系中的公共货币，主要由中央银行负债构成：流通中的通货，以及可在中央银行开户的机构持有的准备金余额。两者都以共同记账单位计价，但在支付与结算系统中履行不同功能。

中央银行发行这些负债，并管理用于银行间结算的准备金账户。一家参与机构向另一家付款时，准备金余额在双方的中央银行账户之间转移。这一转移在机构之间重新分配中央银行负债，并清偿相应的银行间义务。因此，以中央银行货币结算，无须再持有对私人结算中介的额外请求权。

\subsection{私人货币与商业银行}

现代货币体系中的私人货币主要由商业银行存款构成。每笔存款都是存款人的资产和发行银行的负债。作为货币使用的存款以共同记账单位计价，并可通过支付系统转移。尽管不同银行发行的存款采用同一记账单位，每笔存款仍是对特定机构的请求权。

商业银行在扩张资产负债表时创造存款。银行发放贷款时，将贷款记为资产，同时向借款人存款账户记入同额资金，从而创造新的存款负债。存款创造并不要求事先提供等量的中央银行货币。

但存款创造仍受金融和制度约束。随着资产负债表扩张，银行必须维持充足资本和流动性、管理信用风险，并保持履行支付与结算义务的能力。

现代货币体系的一项关键特征是，不同银行发行的存款在法律上属于不同负债，却通常以相同名义价值被接受。这种等价性依赖于连接商业银行存款与中央银行货币、并支持按面值兑换的安排。

\subsection{按面值兑换与货币层级}

不同货币请求权即使是不同发行人的负债，也可能以同一记账单位计价。一家商业银行的人民币 100 元存款，在法律上不同于另一家银行的 100 元存款，也不同于 100 元通货。按面值兑换（par convertibility）意味着这些货币形式通常可以等面值交换，而不按市场决定的折价或溢价交换。

按面值兑换维持名义等价性，却不消除发行人差异。商业银行存款仍是对单家银行的请求权，通货和准备金余额则是中央银行的负债。

这些差异反映为货币层级（hierarchy of money）。商业银行存款可用于家庭和企业的大多数支付，而银行间义务通常以中央银行货币结算。因此，货币层级反映的是发行人和结算地位的差异，而非记账单位的差异。

只要按面值兑换得以维持，这一层级就与货币统一性相容。当对可兑换性的信心减弱时，发行人差异就会变得具有经济意义，货币体系的统一性也可能开始瓦解。

\subsection{维持货币体系的统一性}

按面值兑换使法律上不同的货币负债能够以共同名义价值使用。维持这种货币统一性，取决于私人发行人履行负债的能力，以及结算流动性的可得性。任何一个条件薄弱，都可能破坏货币请求权之间的平价关系。

第一个条件涉及发行银行的财务稳健性。审慎监管限制风险承担并要求充足资本，降低银行无法履行存款负债的可能性。如果银行倒闭，存款保险保护受保障存款，处置安排则为处理倒闭机构提供有序程序。这些机制有助于维持商业银行存款能够按面值兑付的信心。

第二个条件涉及结算流动性。跨行支付产生必须及时清偿的银行间义务。在中央银行持有的准备金余额充当结算资产，流动性工具则可在暂时短缺时提供额外准备金。共同支付规则、互操作性和结算最终性支持可靠的跨行支付，有助于维持按面值兑换。

这些安排共同使法律上不同的货币负债，能够在共同货币体系中作为等价货币发挥作用。

%========================================================
% Summary
%========================================================

\TailModule{summary}{本章小结}

货币不只是被广泛接受的资产。它提供共同记账单位、支持交换，并跨期保存购买力。在现代货币体系中，大多数货币形式也是金融请求权。通货和中央银行准备金是中央银行负债，商业银行存款则是各家银行的负债。它们的有用性依赖于信任、可转移性，以及支持支付和结算的制度安排。

支付不只是传递指令。清算确定参与者之间产生的义务，轧差可以减少待结算金额，结算则通过被接受的结算资产清偿义务。结算的时间和最终性同时影响流动性需求与金融风险。数字支付技术可以改变支付的发起、认证和路由方式，但不会消除这些基本经济功能。

流动性关注按期履行支付与结算义务的能力。可兑换性使货币请求权能够按面值交换，而安全且具有流动性的资产提供额外的流动性和抵押服务。这些服务能够影响资产收益率，并在金融压力时期变得尤为重要。

现代货币体系结合公共货币与私人货币。商业银行通过资产负债表扩张创造存款货币，而中央银行货币提供银行间结算的共同基础。按面值兑换使法律上不同的货币请求权能够以共同名义价值发挥作用。维持这种货币统一性，需要财务稳健的发行人、充足的结算流动性、可靠的支付基础设施，以及能够处理银行困境的制度。

%========================================================
% Key Terms
%========================================================

\TailModule{keyterms}{关键术语}

货币（Money）；记账单位（unit of account）；交换媒介（medium of exchange）；支付系统（payment system）；清算（clearing）；轧差（netting）；结算（settlement）；结算最终性（settlement finality）；流动性（liquidity）；可兑换性（convertibility）；安全资产（safe asset）；便利收益（convenience yield）；中央银行货币（central-bank money）；商业银行货币（commercial-bank money）；按面值兑换（par convertibility）；货币层级。

%========================================================
% Discussion Questions
%========================================================（hierarchy of money

%========================================================
% Discussion Questions
%========================================================）。

\TailModule{questions}{讨论题}

\begin{enumerate}

\item 为什么移动支付应用不一定等同于结算支付所用的资产？数字支付的底层可能有哪些货币请求权？

\item 在延迟净额结算系统中，为什么轧差可以降低流动性要求，却仍使参与者面临结算风险？

\item 商业银行存款虽然是私人机构的负债，为什么仍能充当货币？哪些制度安排支持它按面值被接受？

\item 在金融压力时期，家庭或金融机构为什么可能从银行存款转向通货、其他银行的存款或其他流动资产？这种行为揭示了货币请求权之间的哪些差异？

\end{enumerate}

%========================================================
% Applied Problems
%========================================================

\TailModule{appprobs}{应用题}

\begin{enumerate}

\item 四家银行在清算周期末有以下支付义务：A 欠 B 120；B 欠 C 100；C 欠 D 70；D 欠 A 60；A 还欠 C 30。计算每家银行的付款总额、收款总额和净结算头寸。如果所有义务逐笔全额结算，支付总额是多少？经过完全多边轧差后，哪些银行是净债务人、哪些是净债权人？为完成结算，净债务人总共需要转移多少资金？

\item 银行 X 的一名客户向银行 Y 的客户转账人民币 50,000 元。当银行间义务以准备金余额结算时，列示银行 X、银行 Y 和中央银行的资产负债表变化。再解释：如果两位客户均在银行 X 开户，资产负债表影响有何不同？

\item 两只其他方面可比的一年期证券具有相同预期现金流和相近信用风险。证券 A 的收益率为 3.2\%，可作为中央银行日间信贷的抵押品；证券 B 的收益率为 3.8\%，不具备这一资格。利用式~\eqref{eq:ch2-convenience-yield} 解释 0.6 个百分点的收益率差。要将该差异完全归因于流动性和抵押服务，需要哪些假设？

\end{enumerate}

%========================================================
% Data Application
%========================================================

\TailModule{dataapp}{数据应用}

\textbf{数据中的支付架构：中国，2024 年}

中国人民银行报告显示，2024 年，大额实时支付系统处理业务 3.92 亿笔，金额 8,824.18 万亿元；网联清算平台处理业务约 1.03 万亿笔，金额 520.53 万亿元。\footnote{中国人民银行，《2024 年支付体系运行总体情况》，第三部分“支付系统”。}

\begin{enumerate}

\item 计算两个系统各自的平均每笔交易金额。

\item 比较交易笔数与交易金额。这两个指标揭示了两个系统在经济作用上的哪些差异？


\item 将这些差异与本章讨论的支付系统机制联系起来。为什么大额实时支付系统尤其重视结算流动性和最终性，而处理海量相对小额交易的平台更强调处理能力和较低的单位成本？





\item 使用中国人民银行较新一期支付体系报告重复比较。两个系统的笔数、金额和平均每笔规模如何变化？得出结论之前，检查统计定义或覆盖范围是否发生变化。




\end{enumerate}

%========================================================
% Financial AI Lab
%========================================================

\TailModule{aiprompt}{金融 AI 实验室}

\textbf{实验：追踪支付直至最终结算}

选择一种可以通过可靠来源加以记录的支付方式，例如银行转账、银行卡、二维码支付或数字钱包。请 AI 系统描述从付款人指令到最终结算的完整交易路径。将生成的回答视为待核实的假说，而不是权威描述。

通过五项检查重建交易：识别付款人持有的货币请求权；识别发行该请求权的机构；识别清算安排；识别最终结算使用的资产；识别决定结算最终性的规则或机构。每一步都要求可核实的机构来源，并标注无法确认的陈述。

将核实后的交易路径与 AI 回答比较。识别 AI 在何处将支付界面与底层货币请求权、清算过程、资产负债表转移或结算资产相混淆。最后解释支付过程中哪些部分属于技术，哪些部分依赖金融和制度安排。

%========================================================
% Further Reading
%========================================================

\TailModule{furtherread}{延伸阅读}
\begin{itemize}

\item 中国人民银行（2025），《中国支付体系发展报告（2024）》（\emph{Payment System Report (2024)}）。

\item McLeay, Michael, Amar Radia, and Ryland Thomas (2014),
“现代经济中的货币：导论”（``Money in the Modern Economy: An Introduction''），
\emph{Bank of England Quarterly Bulletin}, 54(1), 4--13.

\item 支付结算体系委员会与国际证券委员会组织技术委员会（2012），《金融市场基础设施原则》（\emph{Principles for Financial Market Infrastructures}），国际清算银行。




\item 国际清算银行（2025），
“下一代货币与金融体系”（``The Next-Generation Monetary and Financial System''），
载《国际清算银行 2025 年度经济报告》（\emph{BIS Annual Economic Report 2025}）第三章。

\end{itemize}

%========================================================
% Tail row in Full Contents
%========================================================

\ChapterTail{summary,keyterms,questions,appprobs,dataapp,aiprompt,furtherread}
